\subsection{Problem S2-SC: explicit adjacent exceptional triangle domains}

For \(0\le c\le1/2\), define
\[
 E_{\rm sc}(c)
 =\frac{c+\sqrt{c^2-8c+4}}2,
 \qquad
 A_{\rm sc}(c)=1-E_{\rm sc}(c).
\]
For \(a\ge0\) and \(u<1\) with \(a+1-u>0\), put
\[
 \vartheta(a,u)=\frac{a}{a+1-u},
\]
and define
\[
 \Psi_{\rm SC}(a,c,u)
 =\max\{a,\vartheta(a,u)\}+E_{\rm sc}(c)-1.
\]

\begin{definition}[First explicit source cell]
\label{def:s2-sc-t3-domain}
For
\[
 0\le\theta\le2-\sqrt3,
\]
define
\[
 x_L(\theta)
 =\frac{1-4\theta+\theta^2}{2(1-2\theta)},
 \qquad
 \rho(\theta)=\frac{1-2\theta}{1-\theta^2}.
\]
Let \(\Omega_{\rm SC}^{T}\) be the set of
\((\theta,x,a,c,u)\in\mathbb R^5\) satisfying
\[
 x_L(\theta)\le x\le\frac12-\theta,
\]
\[
 a=x\rho(\theta),\qquad
 c=x+\theta,\qquad
 u=1-\rho(\theta)+a,
\]
and
\[
 0<a<1,\qquad
 0\le c\le\frac12,\qquad
 \frac12\le u<1.
\]
Every denominator is positive on the displayed \(\theta\)-interval.
\end{definition}

On this cell,
\[
 1-u=(1-x)\rho(\theta),
 \qquad
 \vartheta(a,u)=x.
\]
Thus the variables \((\theta,x)\) replace the earlier translated-form
predicate by an explicit rational domain.

\begin{definition}[Second explicit source cell]
\label{def:s2-sc-vd1-domain}
For \(t\ge1\), put
\[
 \Delta(t)=\sqrt{t^2+t+1}.
\]
Let \(\Omega_{\rm SC}^{V}\) be the set of
\((t,a,b,c,u)\in\mathbb R^5\) satisfying
\[
 0<a<1,\qquad 0<b<1,\qquad 0\le c\le\frac12,
\]
\[
 c=\frac{t(1-b)}{t+1},
 \qquad
 u=\frac{\Delta(t)-a-tb-1}{t},
\]
\[
 0\le c<u<1,
 \qquad
 a+tb\le\Delta(t)-1-\frac t2,
\]
and
\[
 \Delta(t)-a-tb-t
 \le\frac{1-a}{t+1}.
\]
The last inequality is the exact no-positive-interval condition for the
second raw interval; it replaces a geometric type predicate.
\end{definition}

\begin{definition}[Problem S2-SC]
\label{prob:s2-rescuer}
The two optimization statements are
\[
 \boxed{
 \Psi_{\rm SC}(a,c,u)\le0
 \qquad
 ((\theta,x,a,c,u)\in\Omega_{\rm SC}^{T})
 }
\]
and
\[
 \boxed{
 \Psi_{\rm SC}(a,c,u)\le0
 \qquad
 ((t,a,b,c,u)\in\Omega_{\rm SC}^{V}).
 }
\]
Their union is denoted by \(\Omega_{\rm SC}\).
\end{definition}

The first cell is the exact rational parameter range used in the first local
profile proof.  The second cell states the corner equations, midpoint
inequality, positive first interval, and absence of a positive second
interval.  Neither cell contains a translation, support, or placement
predicate.
